\chapter{Analiza problemu i istniejących rozwiązań}
\label{chap:analiza-problemu}

\section{Opis problemu}

Identyfikacja osoby w systemie informatycznym polega na ustaleniu, kto korzysta z systemu lub do kogo należy analizowany zestaw danych. W wielu aplikacjach proces ten opiera się na klasycznych metodach, takich jak login, hasło, karta dostępu, numer identyfikacyjny albo ręczna weryfikacja danych przez operatora. Rozwiązania te są powszechnie stosowane, ale w systemach bezpieczeństwa mogą być niewystarczające.

Hasło może zostać zapomniane, ujawnione lub przechwycone. Karta dostępowa może zostać zgubiona albo przekazana innej osobie. Ręczna weryfikacja danych zależy natomiast od dokładności operatora i może być czasochłonna przy większej liczbie osób. W takich przypadkach sam identyfikator nie zawsze potwierdza, że dana osoba rzeczywiście jest tą, za którą się podaje.

Cechy biometryczne mogą wspierać ten proces, ponieważ są powiązane bezpośrednio z człowiekiem. Twarz i odcisk palca należą do najczęściej wykorzystywanych cech tego typu. Mogą zostać zarejestrowane w postaci obrazu, przetworzone do reprezentacji biometrycznej, a następnie porównane z danymi zapisanymi wcześniej w systemie. Dzięki temu identyfikacja nie opiera się wyłącznie na wiedzy użytkownika lub posiadanym nośniku, ale także na analizie cechy fizycznej.

W projekcie analizowany problem dotyczy identyfikacji osoby na podstawie wcześniej zarejestrowanych danych biometrycznych. System porównuje nową próbkę z danymi zapisanymi w bazie i wskazuje, czy odnaleziono dopasowanie.

\section{Dane biometryczne wykorzystywane w projekcie}

W projekcie wykorzystywane są dwa typy danych biometrycznych:
\begin{itemize}
    \item obraz twarzy,
    \item obraz odcisku palca.
\end{itemize}

Obraz twarzy stanowi dane wejściowe pozwalające na utworzenie reprezentacji osoby na podstawie jej cech wyglądu. Tego typu dane mogą pochodzić z pliku graficznego lub z kamery. Dla poprawnego działania systemu ważne jest, aby twarz była widoczna, odpowiednio oświetlona i możliwa do przetworzenia.

Obraz odcisku palca jest drugim rodzajem danych wejściowych. W tym przypadku system wykorzystuje plik graficzny przedstawiający układ linii papilarnych. Obraz taki służy do utworzenia reprezentacji biometrycznej odcisku, która może zostać porównana z wcześniej zapisanymi danymi.

Oba typy danych nie są traktowane jako zwykłe obrazy porównywane bezpośrednio piksel po pikselu. Ich rola polega na dostarczeniu informacji potrzebnych do utworzenia reprezentacji biometrycznej, która następnie może zostać wykorzystana w procesie identyfikacji. Szczegóły sposobu tworzenia i porównywania tych reprezentacji opisano w dalszym rozdziale dotyczącym metod i algorytmów.

\newpage

\section{Założenia i ograniczenia systemu}

Przyjęto, że system nie pełni roli produkcyjnego systemu bezpieczeństwa, lecz służy do zaprezentowania procesu identyfikacji osób na podstawie danych biometrycznych.

Najważniejsze założenia przyjęte w projekcie są następujące:
\begin{itemize}
    \item system jest obsługiwany przez operatora,
    \item identyfikacja odbywa się na podstawie danych wcześniej zapisanych w bazie,
    \item analizowane dane mają postać obrazów twarzy i odcisków palców,
    \item skuteczność identyfikacji zależy od jakości danych wejściowych,
    \item wynik systemu ma wspierać proces identyfikacji, a nie zastępować pełnej ekspertyzy biometrycznej.
\end{itemize}

Na działanie systemu wpływa jakość obrazu. W przypadku twarzy znaczenie mają między innymi oświetlenie, ostrość zdjęcia, kąt ustawienia głowy oraz zasłonięcie części twarzy. W przypadku odcisku palca istotne są kompletność odcisku, jakość skanu, położenie palca oraz widoczność charakterystycznych elementów linii papilarnych.

Istotnym ograniczeniem jest również brak pełnej detekcji żywotności użytkownika. System analizuje dostarczony obraz, ale nie potwierdza, czy próbka pochodzi od żywej osoby obecnej w momencie identyfikacji.

Możliwym skutkiem ograniczeń jakościowych jest błędna identyfikacja albo odrzucenie poprawnej osoby. Szczegółowa analiza takich zagrożeń oraz odporności systemu została wydzielona do rozdziału poświęconego bezpieczeństwu.

\section{Analiza istniejących rozwiązań}

\subsection{Windows Hello}

Windows Hello jest mechanizmem uwierzytelniania użytkownika w systemie Windows. Może wykorzystywać między innymi rozpoznawanie twarzy, odcisk palca lub kod PIN \cite{windowshello}. Rozwiązanie to jest stosowane głównie do logowania do urządzenia oraz potwierdzania tożsamości użytkownika w środowisku systemu operacyjnego.

Zaletą Windows Hello jest wygoda użytkowania i integracja z systemem Windows. Użytkownik nie musi za każdym razem wpisywać hasła, a proces uwierzytelniania odbywa się szybko i lokalnie na urządzeniu. Ograniczeniem z perspektywy projektu jest zamknięty charakter rozwiązania. Użytkownik nie ma pełnego wglądu w sposób przetwarzania danych biometrycznych, strukturę zapisanych reprezentacji ani szczegóły podejmowania decyzji.

W odniesieniu do projektu IdentityCore Windows Hello stanowi przykład praktycznego wykorzystania biometrii w aplikacjach codziennego użytku. Różnica polega na tym, że IdentityCore nie jest mechanizmem logowania do systemu operacyjnego, lecz aplikacją demonstracyjną pozwalającą obserwować cały proces zarządzania osobami i próbami identyfikacji.

\newpage

\subsection{Face ID}

Face ID jest systemem rozpoznawania twarzy stosowanym w urządzeniach firmy Apple \cite{faceid}. Jego głównym zadaniem jest odblokowywanie urządzenia, autoryzacja operacji oraz potwierdzanie tożsamości użytkownika w ramach ekosystemu Apple. System wykorzystuje biometrię twarzy i jest zintegrowany bezpośrednio ze sprzętem oraz oprogramowaniem urządzenia.

Najważniejszą zaletą Face ID jest wygoda i wysoki poziom integracji z urządzeniem. Proces uwierzytelniania jest szybki, a użytkownik nie musi wykonywać dodatkowych czynności poza spojrzeniem na urządzenie. Ograniczeniem z punktu widzenia dokumentowanego projektu jest to, że Face ID działa jako zamknięty mechanizm systemowy. Nie jest rozwiązaniem służącym do prowadzenia własnego rejestru osób, analizy historii prób ani eksperymentowania z parametrami działania.

Dla projektu IdentityCore Face ID jest punktem odniesienia w zakresie wykorzystania twarzy jako cechy biometrycznej. Projekt nie próbuje odtwarzać tak zaawansowanego systemu sprzętowo-programowego, lecz pokazuje uproszczony proces identyfikacji twarzy w aplikacji desktopowej.

\subsection{Systemy kontroli dostępu z odciskiem palca}

Systemy kontroli dostępu wykorzystujące odciski palców są stosowane między innymi w firmach, laboratoriach, zamkach elektronicznych oraz systemach ewidencji czasu pracy. Ich zadaniem jest sprawdzenie, czy osoba przykładająca palec do czytnika znajduje się w bazie uprawnionych użytkowników \cite{fingerprintaccess}.

Zaletą takich systemów jest prostota użycia oraz bezpośrednie powiązanie próby dostępu z cechą fizyczną użytkownika. Odcisk palca jest wygodny w codziennym użyciu i może być szybko porównany z zapisanymi wzorcami. Ograniczenia dotyczą przede wszystkim jakości odczytu, zabrudzenia palca lub skanera, nieprawidłowego przyłożenia palca oraz ryzyka użycia sztucznie przygotowanego odcisku.

W porównaniu z takimi rozwiązaniami IdentityCore nie jest terminalem kontroli dostępu. Aplikacja nie otwiera drzwi ani nie zarządza uprawnieniami fizycznymi. Wykorzystuje jednak podobny typ danych biometrycznych i pozwala przeanalizować proces identyfikacji odcisku w środowisku projektowym.

\subsection{Komercyjne SDK i systemy biometryczne}

Istnieją również komercyjne zestawy SDK oraz systemy biometryczne umożliwiające budowę aplikacji wykorzystujących rozpoznawanie twarzy, odcisków palców lub innych cech biometrycznych \cite{biometricsdk}. Tego typu rozwiązania oferują gotowe mechanizmy obsługi danych biometrycznych, porównywania próbek, a czasem również integrację z urządzeniami zewnętrznymi.

Ich zaletą jest szeroki zakres funkcji oraz możliwość użycia w profesjonalnych systemach bezpieczeństwa. Ograniczeniem w kontekście projektu jest często zamknięty charakter bibliotek oraz mniejsza przejrzystość działania. Część procesu rozpoznawania pozostaje ukryta wewnątrz gotowego narzędzia.

IdentityCore korzysta z gotowych narzędzi jedynie tam, gdzie jest to uzasadnione zakresem projektu, ale samodzielnie organizuje przepływ danych, rejestr osób, zapis prób i prezentację wyników. Dzięki temu aplikacja pozwala lepiej pokazać zależność między jakością danych, sposobem ich porównania i decyzją systemu.

\newpage

\section{Wyróżniki projektowanego rozwiązania}

IdentityCore różni się od opisanych rozwiązań przede wszystkim tym, że łączy kilka elementów w jednej aplikacji desktopowej. System nie jest wyłącznie mechanizmem logowania ani terminalem kontroli dostępu, lecz aplikacją pozwalającą zarządzać osobami, próbkami biometrycznymi oraz wynikami prób identyfikacji.

Najważniejsze wyróżniki projektowanego rozwiązania to:
\begin{itemize}
    \item połączenie identyfikacji twarzy i odcisku palca w jednej aplikacji,
    \item możliwość zarządzania osobami i próbkami biometrycznymi,
    \item zapis historii prób identyfikacji,
    \item możliwość późniejszej analizy skuteczności działania systemu.
\end{itemize}

Projekt nie konkuruje z komercyjnymi systemami biometrycznymi pod względem poziomu zabezpieczeń ani gotowości produkcyjnej. Jego wartość polega na pokazaniu kompletnego procesu identyfikacji w aplikacji, która łączy część biometryczną, użytkową i analityczną.
